% ======================================================================
% Ontologie der Kontinua -- Core 1.3
% Cross-K-Modul: crossk_k6_k7.tex
% Cross-Level-Relationen zwischen K6 und K7
% ======================================================================

\subsubsection{K6–K7-Querverbindung}
\label{sec:crossk-k6-k7}

Der Übergang $K_6 \rightarrow K_7$ markiert die Entstehung des sozialen
Kontinuums aus dem kognitiven Kontinuum. Während $K_6$ konzeptionelle Strukturen,
Bedeutung und interne Modelle unterstützt, führt $K_7$ ein:
\begin{itemize}
  \item soziale Rollen und Normen als neue Achsen,
  \item Kommunikationsflüsse und geteilte symbolische Strukturen,
  \item kollektive Schwellen für Koordination und Erwartung,
  \item institutionelle Zyklen und stabilisierende Rückkopplungsschleifen,
  \item Muster von Kooperation, Konflikt und Normdurchsetzung,
  \item einen gemeinsamen Raum $\Omega(K_7)$ sozialer gültiger Konfigurationen.
\end{itemize}

Dieser Übergang entspricht dem Auftreten kollektiver Organisation, die nicht auf
individuelle Kognition reduziert werden kann.

% ----------------------------------------------------------------------
\subsubsection{1. Stufen und ihre Rollen}

\paragraph{K6 (kognitives Kontinuum).}
$K_6$ liefert:
\begin{itemize}
  \item konzeptionelle Achsen $A^{c}$,
  \item kognitive Potentiale $P^{c}$,
  \item Schwellen $\Theta^{c}$ für Kohärenz und Widerspruch,
  \item Flüsse von Aufmerksamkeit, Gedächtnis, Interpretation $J^{c}$,
  \item Zyklen von Verstehen und Lernen,
  \item Kohärenzmaß $k_6$.
\end{itemize}

Allerdings fehlen in $K_6$:
\begin{itemize}
  \item Inter-Agentur-Koordinationsregeln,
  \item gemeinsame symbolische Normen,
  \item populationsbezogene Stabilitätsgrenzen,
  \item institutionelle Persistenz.
\end{itemize}

\paragraph{K7 (soziales Kontinuum).}
$K_7$ führt ein:
\begin{itemize}
  \item Achsen $A_{\mathrm{soc}}$ (Rollen, Normen, Status, Identität),
  \item Potentiale $P_{\mathrm{soc}}$ (Vertrauen, Autorität, Legitimität),
  \item soziale Schwellen $\Theta_{\mathrm{soc}}$ (Koordination, Kohäsion,
        Durchsetzung, Kollaps),
  \item Flüsse $J_{\mathrm{comm}}$ (Kommunikation, Signalisierung, Imitation),
  \item Zyklen $C_{\mathrm{inst}}$ (Regel–Aktion–Konsequenz–Korrektur),
  \item ein globales Maß $k_7$ für soziale Stabilität.
\end{itemize}

$K_7$ ist die erste Stufe, in der Erwartungen und Normen Verhalten über mehrere
Akteure regulieren.

% ----------------------------------------------------------------------
\subsubsection{2. Vererbte und neue Achsen sowie Schwellen}

\paragraph{Vererbung.}
Konzeptionelle Unterscheidungen aus $K_6$ liefern das Substrat für geteilte Symbole:
\[
A^{c}(K_6) \rightarrow A_{\mathrm{soc}}(K_7).
\]
Interne Modelle werden zur Grundlage geteilter Narrative.

\paragraph{Neue Achsen.}
$K_7$ führt ein:
\begin{itemize}
  \item $A_{\mathrm{role}}$ — Verhaltenserwartungen,
  \item $A_{\mathrm{norm}}$ — erlaubte/unerlaubte Handlungen,
  \item $A_{\mathrm{status}}$ — soziale Differenzierung,
  \item $A_{\mathrm{collective}}$ — Gruppenidentitätsstrukturen.
\end{itemize}

Diese Achsen existieren nicht in $K_6$.

\paragraph{Schwellen.}
$\Theta_{\mathrm{soc}}$ enthält:
\begin{itemize}
  \item Kohäsionsschwellen (minimales Vertrauen oder Kommunikationsdichte),
  \item Widerspruchsschwellwerte (Inkompatibilität von Normen),
  \item Durchsetzungsgrenzen (Stabilität von Sanktionen),
  \item Kollapsschwellen (Versagen institutioneller Zyklen).
\end{itemize}

Verletzung von $\Theta^{c}$ lässt Kohärenz auf kognitiver Ebene zusammenbrechen
und verhindert das Auftreten von $K_7$.

% ----------------------------------------------------------------------
\subsubsection{3. Cross-Level-Flüsse, Zyklen und Spannungen}

\paragraph{Flüsse.}
Soziale Flüsse $J_{\mathrm{comm}}$ erweitern kognitive Flüsse:
\[
J_{\mathrm{comm}} = (J_{\mathrm{signal}},\ J_{\mathrm{language}},\
                     J_{\mathrm{imitation}},\ J_{\mathrm{coord}}).
\]
Sie verteilen Bedeutung über Akteure und stabilisieren kollektive Zustände.

\paragraph{Zyklen.}
$K_7$ ist durch institutionelle Zyklen gekennzeichnet:
\[
C_{\mathrm{inst}}:\ \text{Regel} \rightarrow \text{Aktion}
\rightarrow \text{Konsequenz} \rightarrow \text{Korrektur}
\rightarrow \text{Regel}.
\]
Eine Gesellschaft existiert nur, wenn mindestens ein solcher stabilisierender
Zyklus intakt bleibt.

\paragraph{Spannung.}
Cross-Level-Spannung $T_{6,7}$ misst Inkongruenz zwischen persönlichen
konzeptionellen Modellen und kollektiven Strukturen:
\[
T_{6,7} = f(A_{\mathrm{soc}} - A^{c},\ P_{\mathrm{soc}} - P^{c},\
          J_{\mathrm{comm}} - J^{c},\ \Theta_{\mathrm{soc}} - \Theta^{c}).
\]
Überschreitet $T_{6,7}$ die dimensionsbezogene Schwelle, scheitert
kollektive Organisation und das System kollabiert zur individuellen Kognition.

% ----------------------------------------------------------------------
\subsubsection{4. Geburts- und Todesbedingungen über die Stufen}

\paragraph{Geburt von \texorpdfstring{$K_7$}{K_7}.}
$K_7$ entsteht, wenn:
\[
T_6 > \Theta_{6,\mathrm{dim}}
\quad\text{und}\quad
A_{\mathrm{soc}} \in M_7 \setminus A(K_6).
\]
Das entspricht der Bildung stabiler sozialer Normen und Rollen.

Erweiterung des Zustandsraums:
\[
\Omega(K_7) = \Omega(K_6) \cup \Delta\Omega_{\mathrm{soc}}.
\]

\paragraph{Todesausbreitung.}
Wenn $\Theta^{c}$ versagt (kognitiver Kollaps), wird $K_7$ unmöglich.

Wenn $\Theta_{\mathrm{soc}}$ versagt (Verlust von Kohäsion, Zusammenbruch von
Zyklen), kollabiert das System zu individueller Kognition:
\[
\Omega(K_7) \rightarrow \Omega(K_6).
\]

$K_6$ bleibt stabil, solange die eigenen Schwellen nicht verletzt sind.

% ----------------------------------------------------------------------
\subsubsection{5. Konzeptionelle Beispiele}

\paragraph{Entstehung von Normen.}
Wiederholte kognitive Interpretationen über Akteure hinweg stabilisieren sich zu
kollektiven Normen und bilden neue Achsen in $A_{\mathrm{soc}}$.

\paragraph{Institutionelle Stabilität.}
Wenn ein institutioneller Zyklus $C_{\mathrm{inst}}$ schließt mit minimaler Drift:
\[
\oint_C dA_{\mathrm{soc}} \approx 0,
\]
bildet sich ein stabiles soziales Subsystem.

\paragraph{Zusammenbruch.}
Sinkt die Kommunikationsdichte unter eine Kohäsionsschwelle:
\[
J_{\mathrm{comm}} < \Theta_{\mathrm{soc,coh}},
\]
löst sich die kollektive Struktur auf.

Der K6→K7-Übergang lässt sich so in komprimierter Form als Entstehen geteilter
Normen, Kommunikationsflüsse und institutioneller Zyklen aus individueller Kognition
zusammenfassen.
